\SourceLeaf{436}%
\Needspace{7\baselineskip}\begin{center}13.\end{center}
\noindent\hspace*{1.2em}In our representation, every point of a curved line on the curved surface already has two corresponding points on the sphere. The point $\lambda$ represents the direction of the line element, and the point $L$ represents the direction of the normal to the surface. They are evidently $90^{0}$ apart. In our earlier investigation (Art. $9$), where [we] assumed the curved line to lie in a single plane, we had \textls[100]{two} further points on the sphere: $\mathfrak{L}$, representing the direction of the normal to the plane, and $\lambda'$, representing the direction of the normal to the curve element in the plane. In that case $\mathfrak{L}$ was fixed and $\lambda$, $\lambda'$ always lay on a single great circle whose pole was $\mathfrak{L}$. In generalizing the investigation, we shall retain the notations $\mathfrak{L}$, $\lambda'$, but must understand their meaning from a more general point of view. As the curve $s$ is traced, the points $L,\lambda$ on the auxiliary sphere also trace curves that, generally speaking, are no longer great circles. Draw a radius of the auxiliary sphere parallel to the element of the second curve, ending at the point $\lambda'$. We choose the opposite point instead if that point is more than $90^{0}$ from $L$. In the first case we take the element at $\lambda$ to be positive, and in the other case negative. Finally, let $\mathfrak{L}$ be the point on the auxiliary sphere separated from both $\lambda$ and $\lambda'$ by $90^{0}$, with $\lambda,\lambda',\mathfrak{L}$ in the same order as $(1),(2),(3)$.

The coordinates of the four points on the auxiliary sphere relative to its centre are
\[
\begin{array}{cccc}
L&X&Y&Z\\
\lambda&\xi&\eta&\zeta\\
\lambda'&\xi'&\eta'&\zeta'\\
\mathfrak{L}&\alpha&\beta&\gamma.
\end{array}
\]
Each of these $4$ points therefore traces a line on the auxiliary sphere itself. We shall denote their elements by $dL,d\lambda,d\lambda',d\mathfrak{L}$. We thus have
